\documentclass[10.5pt,a4paper]{article}
\usepackage[a4paper,margin=22mm]{geometry}
\usepackage{amsmath,amssymb,amsthm,mathtools}
\usepackage{fontspec}
\usepackage{xetexko}
\setmainhangulfont{NanumMyeongjo}
\setsanshangulfont{NanumGothic}
\setmonohangulfont{NanumGothicCoding}
\usepackage{booktabs,longtable,array,graphicx,enumitem,xcolor,float}
\usepackage[colorlinks=true,linkcolor=blue!50!black,urlcolor=blue!50!black,citecolor=blue!50!black]{hyperref}
\setlist{itemsep=1pt,topsep=3pt}
\linespread{1.12}
\setlength{\parskip}{3pt}
\setlength{\parindent}{0pt}

\theoremstyle{plain}
\newtheorem{theorem}{정리}[section]
\newtheorem{proposition}[theorem]{명제}
\newtheorem{lemma}[theorem]{보조정리}
\newtheorem{corollary}[theorem]{따름정리}
\newtheorem{conjecture}[theorem]{추측}
\theoremstyle{definition}
\newtheorem{definition}[theorem]{정의}
\newtheorem{example}[theorem]{예}
\newtheorem{observation}[theorem]{수치관찰}
\newtheorem{heuristic}[theorem]{휴리스틱}
\theoremstyle{remark}
\newtheorem{remark}[theorem]{주의}


\usepackage{newunicodechar}
\newunicodechar{≈}{\ensuremath{\approx}}
\newunicodechar{⇒}{\ensuremath{\Rightarrow}}
\newunicodechar{→}{\ensuremath{\to}}
\newunicodechar{↔}{\ensuremath{\leftrightarrow}}
\newunicodechar{≤}{\ensuremath{\le}}
\newunicodechar{≥}{\ensuremath{\ge}}
\newunicodechar{×}{\ensuremath{\times}}
\newunicodechar{∈}{\ensuremath{\in}}
\newunicodechar{⊂}{\ensuremath{\subset}}
\newunicodechar{≠}{\ensuremath{\ne}}
\newunicodechar{∞}{\ensuremath{\infty}}
\newunicodechar{√}{\ensuremath{\surd}}
\newunicodechar{≪}{\ensuremath{\ll}}
\newunicodechar{≫}{\ensuremath{\gg}}
\newunicodechar{∼}{\ensuremath{\sim}}
\newunicodechar{⊕}{\ensuremath{\oplus}}
\newunicodechar{⊗}{\ensuremath{\otimes}}
\newunicodechar{∅}{\ensuremath{\emptyset}}
\newunicodechar{∩}{\ensuremath{\cap}}
\newunicodechar{∪}{\ensuremath{\cup}}
\newunicodechar{±}{\ensuremath{\pm}}
\providecommand{\tightlist}{\setlength{\itemsep}{0pt}\setlength{\parskip}{0pt}}
\usepackage{calc}
\usepackage{url}
\usepackage{fvextra}
\AtBeginDocument{\renewcommand{\figurename}{그림}\renewcommand{\contentsname}{차례}\renewcommand{\abstractname}{초록}}
\newfontfamily\notesfont{DejaVu Sans Mono}[Scale=0.82]

\newcommand{\R}{\mathbb{R}}
\newcommand{\Z}{\mathbb{Z}}
\newcommand{\Q}{\mathbb{Q}}
\newcommand{\C}{\mathbb{C}}
\newcommand{\HH}{\mathbb{H}}
\newcommand{\E}{\mathbb{E}}
\renewcommand{\P}{\mathbb{P}}
\newcommand{\Sym}{\mathrm{Sym}}
\newcommand{\Tr}{\mathrm{Tr}}
\newcommand{\Vol}{\mathrm{Vol}}
\newcommand{\covol}{\mathrm{covol}}
\newcommand{\rd}{\mathrm{rd}}
\newcommand{\Herm}{\mathrm{Herm}}
\newcommand{\tagP}{\textcolor{green!45!black}{\textbf{[증명됨]}}}
\newcommand{\tagPC}{\textcolor{green!45!black}{\textbf{[증명됨(계산기)]}}}
\newcommand{\tagC}{\textcolor{orange!80!black}{\textbf{[조건부]}}}
\newcommand{\tagJ}{\textcolor{purple}{\textbf{[추측]}}}
\newcommand{\tagN}{\textcolor{blue!70!black}{\textbf{[수치관찰]}}}
\newcommand{\tagH}{\textcolor{gray!70!black}{\textbf{[휴리스틱]}}}
\newcommand{\tagL}{\textcolor{teal!70!black}{\textbf{[문헌]}}}
\newcommand{\tagF}{\textcolor{red!70!black}{\textbf{[실패/반증]}}}
\newcommand{\tagM}{\textcolor{green!45!black}{\textbf{[모형에서 증명]}}}
\newcommand{\tagPL}{\textcolor{green!45!black}{\textbf{[증명됨·인용 정리 전제]}}}
\providecommand{\Hom}{\operatorname{Hom}}
